% Work/MAT – přepracovaná 21. hodina: 20min průběžný test z množin + 25min úvod GON.
% Reviduje původní 21_MNZ_PrubeznyTestARozbor.tex, zachovává klasifikovaný test.
% Zadání testu se řídí školní schválenou verzí, zde se nevytváří nový test.
% Učebnice: ucebnice/04_Goniometrie_a_trigonometrie/Goniometrické funkce.tex;
%               ucebnice/03_Planimetrie/Planimetrie.tex.
% Realisticky.cz: 010119 Podobnosti trojúhelníků, goniometrické funkce;
%                 040201 Goniometrické funkce ostrého úhlu (indexovaný přehled).
% Příklady, kresby a formulace jsou vlastní.
\documentclass[11pt,a4paper]{article}
% Spolecny styl pro ucitelske pripravy MAT (Work/MAT)
% Verze 1.1, 2026-10-23
\NeedsTeXFormat{LaTeX2e}
\ProvidesPackage{priprava}[2026/10/23 v1.1 Sablona ucitelskych priprav MAT]

% Zakladni typografie a cestina
\RequirePackage[T1]{fontenc}
\RequirePackage[utf8]{inputenc}
\RequirePackage[czech]{babel}
\RequirePackage{lmodern}
\RequirePackage{microtype}
\RequirePackage{amsmath,amssymb}
\RequirePackage{xcolor}
\RequirePackage{enumitem}
\RequirePackage[most]{tcolorbox}
\RequirePackage{tabularx,booktabs,array}
\RequirePackage{fancyhdr}
\RequirePackage{needspace}
\RequirePackage{geometry}
\RequirePackage[hidelinks]{hyperref}

\geometry{a4paper,top=17mm,bottom=18mm,left=23mm,right=23mm}
\definecolor{notegray}{gray}{0.39}
\setlength{\parindent}{0pt}
\setlength{\parskip}{2.5pt plus 1pt minus 0.5pt}
\setlist[itemize]{topsep=2pt,itemsep=1pt,parsep=0pt,leftmargin=1.45em}
\setlist[enumerate]{topsep=2pt,itemsep=2pt,parsep=0pt,leftmargin=1.65em}
\setcounter{secnumdepth}{0}
\renewcommand{\arraystretch}{1.12}

% Zahlavi: vlevo rocnik | globalni tema, vpravo cislo: nazev hodiny
% Zapati zustava zcela prazdne.
\pagestyle{fancy}
\fancyhf{}
\renewcommand{\headrulewidth}{0.35pt}
\renewcommand{\footrulewidth}{0pt}
\setlength{\headheight}{14pt}
\newcommand{\nastavPripravu}[4]{%
  \fancyhead[L]{\small #1\enspace\textbar\enspace #2}%
  \fancyhead[R]{\small #3: #4}%
}

% Minimalisticke cernobile boxy. Kazda definice zvlast.
\tcbset{commonlesson/.style={enhanced,breakable=false,colback=white,colframe=black,
  colbacktitle=white,coltitle=black,boxrule=0.45pt,
  arc=0.8mm,outer arc=0.8mm,left=3mm,right=3mm,top=1.6mm,bottom=1.6mm,
  toptitle=1mm,bottomtitle=0.7mm,fonttitle=\bfseries,
  attach boxed title to top left={xshift=2.5mm,yshift=-1.15mm},
  boxed title style={colback=white,colframe=white,boxrule=0pt,left=1mm,right=1mm}}}
\newtcolorbox{zadaniBox}[1][]{commonlesson,title={Zadání},#1}
\newtcolorbox{prikladBox}[1][]{commonlesson,title={Příklad},#1}
\newtcolorbox{definiceBox}[1][]{commonlesson,title={Definice},#1}

% Spolecny box pro zadani a jeho reseni, bez mezery mezi castmi.
% Pouziti:
%   \begin{prikladReseniBox}
%     Zadani ulohy.
%     \tcblower
%     \begin{align*} ... \end{align*}
%   \end{prikladReseniBox}
% Spodni cast ma jemny sedy podklad a delici linku; lze v ni sazet
% vice radku matematiky. Obe casti maji jen jeden spolecny ramecek.
\tcbset{solutionlesson/.style={
  commonlesson,
  enhanced jigsaw,
  breakable,
  colbacklower=black!3,
  segmentation style={solid,draw=black!25,line width=0.35pt},
  fontlower=\small\color{notegray},
  before lower={\noindent\textbf{Řešení:}\enspace},
  leftlower=3mm,rightlower=3mm
}}
\newtcolorbox{prikladReseniBox}[1][]{solutionlesson,title={Příklad},#1}
\newtcolorbox{zadaniReseniBox}[1][]{solutionlesson,title={Zadání},#1}

% Pomucky pro autora pripravy: sedive poznamky; \reseni pro starsi pripravy
\newcommand{\poznamka}[1]{\par\begingroup\small\color{notegray}\itshape
  \textnormal{\textbf{Poznámka pro učitele.}}\enspace #1\par\endgroup}
\newcommand{\reseni}[1]{\par\begingroup\small\color{notegray}
  \textbf{Očekávané řešení:}\enspace #1\par\endgroup}
\newcommand{\blok}[1]{\needspace{5\baselineskip}\subsection*{#1}}
\newcommand{\vyrok}[1]{\uv{#1}}

\usepackage{tikz}
\usetikzlibrary{angles,quotes,babel}
\nastavPripravu{1. ročník}{Množiny / goniometrie}{21}{Test a úvod do sinu a kosinu}
\begin{document}
\blok{První část: průběžný test z množin (20 minut)}
\textbf{0.--20. minuta:} písemně ověřujeme látku hodin 12--20,
včetně intervalů připomenutých v hodině 8. Schválené žákovské zadání
je samostatný soubor a není součástí této učitelské přípravy.

\begin{zadaniBox}[title={Rámec testu – pro učitele}]
Vybírej úlohy na prvky a podmnožiny, základní operace,
Vennovy diagramy a počty prvků, operace s intervaly a přiměřené
slovní úlohy. Nezkoušej goniometrii: začíná až po testu.
Délka čistého řešení: \textbf{20 minut}, podle školních pravidel.
\end{zadaniBox}
\poznamka{Včas připrav a rozdej existující verze testu. Během testu se
nevyučuje nové učivo. V 20. minutě práce vyber. Rozbor výsledků testu
proběhne později podle skutečných chyb; neukrajuj zde čas určený
úvodní goniometrii.}

\blok{Druhá část: jak vůbec můžeme popsat úhel? (25 minut)}
\textbf{20.--23. minuta (diskuse):} Setkali jste se už se slovy
\emph{sinus, kosinus, tangens, kotangens}? Kde? K čemu by
mohly být užitečné? Uměli byste určit výšku budovy, když známe
vzdálenost a úhel pohledu?

\poznamka{Není cílem hned vysvětlit všechny čtyři funkce.
Dnes pouze sinus a kosinus ostrého úhlu v pravoúhlém trojúhelníku,
několik přesných hodnot a orientace v tabulce. Pythagorovu větu a práci
s délkami odvěsen připomeneme v hodině 22.}

\begin{definiceBox}
\textbf{Sinus ostrého úhlu $\alpha$} v pravoúhlém trojúhelníku
je poměr délky odvěsny \emph{protilehlé} k $\alpha$ a délky přepony:
\[
\sin\alpha=\frac{\text{protilehlá odvěsna}}{\text{přepona}}.
\]
\end{definiceBox}
\begin{definiceBox}
\textbf{Kosinus ostrého úhlu $\alpha$} v pravoúhlém trojúhelníku
je poměr délky odvěsny \emph{přilehlé} k $\alpha$ a délky přepony:
\[
\cos\alpha=\frac{\text{přilehlá odvěsna}}{\text{přepona}}.
\]
\end{definiceBox}

\begin{center}
\begin{tikzpicture}[scale=.77]
\coordinate (A) at (0,0);\coordinate (B) at (5,0);\coordinate (C) at (5,3.2);
\draw[thick] (A)--node[below] {přilehlá odvěsna} (B)
 --node[right] {protilehlá odvěsna} (C)
 --node[above left] {přepona} cycle;
\draw (4.65,0)--(4.65,.35)--(5,.35);
\pic[draw,angle radius=8mm,"$\alpha$",angle eccentricity=1.5]{angle=B--A--C};
\node[below left] at (A) {$A$};\node[below right] at (B) {$B$};
\node[above right] at (C) {$C$};
\end{tikzpicture}
\end{center}
\textbf{Proč jsou to funkce úhlu?} Podobné pravoúhlé trojúhelníky
se stejným ostrým úhlem mají stejné poměry odpovídajících stran.
Po zvětšení trojúhelníku se sinus ani kosinus daného úhlu nezmění.

\blok{Základní úhly: první setkání s přesnými hodnotami}
\textbf{30.--38. minuta:} Hodnoty pro $30^\circ$ a $60^\circ$ vycházejí
z poloviny rovnostranného trojúhelníku; pro $45^\circ$ z
rovnoramenného pravoúhlého trojúhelníku. Dnes přijmi výsledky jako
přehled a ověř na jednom obrázku, která odvěsna je protilehlá.

\begin{center}
\renewcommand{\arraystretch}{1.42}
\begin{tabular}{@{}lccccc@{}}
\toprule
Úhel $\alpha$ & $0^\circ$ & $30^\circ$ & $45^\circ$ & $60^\circ$ & $90^\circ$ \\
\midrule
$\sin\alpha$ & $0$ & $\frac12$ & $\frac{\sqrt2}{2}$ & $\frac{\sqrt3}{2}$ & $1$ \\
$\cos\alpha$ & $1$ & $\frac{\sqrt3}{2}$ & $\frac{\sqrt2}{2}$ & $\frac12$ & $0$ \\
\bottomrule
\end{tabular}
\end{center}
\textbf{Důležitá přesnost:} poměrové definice výše platí
pro $0^\circ<\alpha<90^\circ$. Hodnoty pro $0^\circ$ a $90^\circ$
jsou krajní hodnoty funkcí, nikoli hodnoty definované pomocí
skutečného pravoúhlého trojúhelníku s takovým ostrým úhlem.
Později je vysvětlíme geometricky pomocí jednotkové kružnice.

\begin{prikladReseniBox}[title={P1: jak číst tabulku},breakable=false]
Bez kalkulačky urči $\sin30^\circ$, $\cos60^\circ$, $\sin45^\circ$
a $\cos0^\circ$. Co je společné pro $\sin30^\circ$ a $\cos60^\circ$?
\tcblower
$\sin30^\circ=\frac12$, $\cos60^\circ=\frac12$,
$\sin45^\circ=\frac{\sqrt2}{2}$, $\cos0^\circ=1$.
První dva výsledky jsou stejné, protože úhly $30^\circ$ a $60^\circ$
jsou v pravoúhlém trojúhelníku doplňkové do $90^\circ$.
\end{prikladReseniBox}

\blok{Rychlá samostatná práce a závěr (38.--45. minuta)}
\begin{zadaniReseniBox}[title={S1: doplň bez kalkulačky},breakable=false]
\begin{enumerate}[label=\alph*)]
\item $\sin60^\circ=\;?$ a $\cos30^\circ=\;?$.
\item $\sin0^\circ=\;?$ a $\cos90^\circ=\;?$.
\item Platí $\sin45^\circ=\cos45^\circ$?
\item Co označuje \emph{protilehlá} odvěsna? Záleží na zvoleném úhlu?
\end{enumerate}
\tcblower
\begin{enumerate}[label=\alph*)]
\item Obě $\frac{\sqrt3}{2}$.
\item Obě $0$.
\item Ano, obě jsou $\frac{\sqrt2}{2}$.
\item Leží proti zvolenému ostrému úhlu; při změně zvoleného úhlu
se prohodí, která odvěsna je protilehlá a která přilehlá.
\end{enumerate}
\end{zadaniReseniBox}
\poznamka{Čas: test 20 min. Následně 3 min motivační rozhovor,
7 min definice a trojúhelník, 8 min tabulka a P1, 5 min samostatná S1,
2 min závěr. Celkem 45 min. Nedělej teď výpočet obecné délky ze sinu,
tangens ani práci s radiány.}

\blok{Krátká domácí práce}
\begin{zadaniReseniBox}[title={D1: dvě rychlé otázky},breakable=false]
\textbf{a)} Vysvětli, proč $\sin30^\circ=\cos60^\circ$.
\quad \textbf{b)} Urči $\sin90^\circ+\cos90^\circ$.
\tcblower
\textbf{a)} Úhly se doplňují do $90^\circ$, proto jde
v trojúhelníku o stejnou odvěsnu vůči stejné přeponě;
obě hodnoty jsou $\frac12$.
\quad \textbf{b)} $1+0=1$.
\end{zadaniReseniBox}
\end{document}
